% Complete Section 7.8 and Box 7.1, original PDF 103-106 / printed 93-96.
\endgroup
\begingroup\emergencystretch=2em
% block: historical-antecedents-heading page:103 type:section
\subsection{{Historical antecedents}}\label{sec:historical-antecedents}

% block: finite-element-families-unification page:103 type:ordinary-paragraph
{We have defined the simplicial finite element spaces $\mathcal P_r\Lambda^k(\mathcal T_h)$ and $\mathcal P_r^-\Lambda^k(\mathcal T_h)$ in a systematic and unified fashion in any dimension $n$ for every $0\leq k\leq n$ and $r\geq1$. In one, two, and three dimensions these spaces of finite element differential forms may be identified via proxies with spaces of scalar- and vector-valued finite element spaces. In this way, they unify a broad range of finite element methods, developed relatively independently by many researchers over a period of many decades.}

% block: classical-finite-element-antecedents pages:103,104 type:ordinary-paragraph
{In any number of dimensions, the spaces $\mathcal P_r\Lambda^0$ and $\mathcal P_r^-\Lambda^0$ coincide with each other and identify with the most classical finite element space, the Lagrange space of continuous piecewise polynomials of degree at most $r$, going back in the case $r=1$ to Courant \BookCite{36} in 1943. These are the most natural finite elements to use to discretize the space $H^1$. The spaces $\mathcal P_r\Lambda^n$ and $\mathcal P_{r+1}^-\Lambda^n$ (note the change of index) also coincide with each other, and they identify with the larger space of all piecewise polynomials of degree at most $r$, whether continuous or not, the natural finite element discretization of $L^2$. Such spaces are used in finite element calculations in the discontinuous Galerkin method, initiated by Reed and Hill \BookCite{65} in 1973. The spaces $\mathcal P_r^-\Lambda^1$ in two dimensions identify with the famous space of Raviart and Thomas \BookCite{64}, introduced in 1975, to discretize $H(\operatorname{div})$. That space was generalized in two different ways to three dimensions in 1980 by Nédélec \BookCite{59}, once as a discretization of $H(\operatorname{div})$, the other as a discretization of $H(\operatorname{curl})$. These correspond to the spaces $\mathcal P_r^-\Lambda^1$ and $\mathcal P_r^-\Lambda^2$ in three dimensions, respectively. Finally, the two-dimensional space $\mathcal P_r\Lambda^1$ identifies with the Brezzi–Douglas–Marini space from 1985 and the three-dimensional spaces $\mathcal P_r\Lambda^1$ and $\mathcal P_r\Lambda^2$ with their generalization by Nédélec in 1986, referred to as the Nédélec edge and face elements of the second kind. Thus, in $n\leq3$ dimensions, all these spaces have appeared in the finite element literature. The FEEC treatment revealed the unity among these spaces, which had arisen in diverse contexts. In some cases, the most natural DOFs for the spaces were not known before the development of FEEC.}

% block: geometric-and-feec-historical-antecedents page:104 type:ordinary-paragraph
{It is also worth noting that some of these spaces of piecewise polynomial differential forms appeared independently in the context of geometry. The spaces $\mathcal P_1^-\Lambda^k$ are highly geometric, having one DOF per $k$-simplex in the triangulation, and so being isomorphic to the space of simplicial $k$-cochains. Indeed, these spaces were constructed by Whitney in his 1957 book \BookCite{70}, decades before their appearance as mixed finite elements. The connection to mixed finite elements was pointed out much later yet, by Bossavit \BookCite{20} in 1988. Further, the complete polynomial spaces $\mathcal P_r\Lambda^k$ associated to a triangulation make an appearance in the topological work of Sullivan \BookCite{67}, \BookCite{68} in the 1970s, a decade before their reinvention as Brezzi–Douglas–Marini elements. In a largely overlooked conference proceedings from the early 1980s \BookCite{17}, Baker called these Sullivan–Whitney forms and he analyzed the convergence of a discretization of an eigenvalue problem for the Hodge Laplacian using them. A first unified treatment of the $\mathcal P_r^-\Lambda^k$ spaces was made in a seminal paper of Hiptmair \BookCite{50} in 1999. In the 2006 paper of Arnold, Falk, and Winther \BookCite{11} in which the term \emph{finite element exterior calculus} first appeared, the Koszul complex was first applied to finite elements, allowing a major simplification and a unified treatment of both the $\mathcal P_r^-\Lambda^k$ and $\mathcal P_r\Lambda^k$ spaces.}

% block: cubical-finite-element-antecedents page:104 type:ordinary-paragraph
{For the cubic finite element spaces, we briefly described the two families $\mathcal Q_r^-\Lambda^k(\mathcal T_h)$ and $\mathcal S_r\Lambda^k(\mathcal T_h)$. The former can be considered a tensor product analogue of the $\mathcal P_r^-\Lambda^k(\mathcal T_h)$ spaces and had been derived alongside the corresponding simplicial spaces. For example, in two dimensions the $\mathcal Q_r^-\Lambda^1$ spaces are the quadrilateral Raviart–Thomas spaces from \BookCite{64}. The $\mathcal S_r\Lambda^k$ family is different in that it was largely absent from the literature and its discovery was inspired by FEEC. The elements for 0-forms, ordinary $H^1$ finite elements, are well-known in two dimensions. They are the popular serendipity finite elements. The construction of $\mathcal S_r\Lambda^0$ gives the systematic generalization of these elements to higher dimensions. The $\mathcal S_r$ spaces for 1-forms and 2-forms in three and more dimensions were likewise not previously known.}

% block: finite-element-forms-summary-heading page:105 type:ordinary-paragraph
\Needspace{32\baselineskip}
\refstepcounter{booksummary}\label{box:finite-element-forms-summary}
\begin{tcolorbox}[breakable,frame hidden,colback=KuangNei]
\textbf{Box~\ref{box:finite-element-forms-summary}. {Summary of Chapter~\ref{chap:finite-element-differential-forms} on finite element differential forms.}}

% block: finite-element-forms-summary-spaces page:105 type:ordinary-paragraph
{Give an $n$-dimension simplicial triangulation $\mathcal T_h$ for any $n\geq1$, any form degree $0\leq k\leq n$, and any polynomial degree $r\geq1$, there are two canonical spaces of finite element differential $k$-forms consisting of piecewise polynomials of degree $r$ or less, the complete polynomial space $\mathcal P_r\Lambda^k(\mathcal T_h)$ and the trimmed polynomial space $\mathcal P_r^-\Lambda^k(\mathcal T_h)$. Like any finite element spaces, these can be specified by giving their shape functions and DOFs and showing that these are unisolvent.}

% block: finite-element-forms-summary-trimming page:105 type:ordinary-paragraph
{As the name suggests, the shape functions for the complete polynomial space consist of all $k$-forms with coefficients which are polynomials of degree at most $r$. The trimmed polynomial space restricts the coefficient polynomials to a subspace of all polynomials of degree at most $r$ but which contains at least all polynomials of degree at most $r-1$. It is precisely defined in terms of the Koszul complex. In the case of 0-forms, the trimmed space coincides with the complete polynomial space of the same degree, while for the case of $n$-forms, it coincides with the complete polynomial space of degree $r-1$. In all other cases, the trimmed space is not a complete polynomial space. In the lowest order case $r=1$, the trimmed spaces are the spaces of Whitney forms.}

% block: finite-element-forms-summary-complexes page:105 type:ordinary-paragraph
{The complete polynomial spaces with $r$ decreasing as $k$ increases combine to form a subcomplex of the $L^2$ de Rham complex with bounded cochain projections. The trimmed polynomial spaces similarly form a de Rham subcomplex, but with degree $r$ constant.}

% block: finite-element-forms-summary-methods page:105 type:ordinary-paragraph
{Two such spaces, one for $k$-forms and one for $(k-1)$-forms, either both complete, both trimmed, or one of each, can be used to solve the $k$-form Hodge Laplacian, as indicated in \eqref{eq:four-galerkin-form-pairs}.}

% block: finite-element-forms-summary-cubical page:105 type:ordinary-paragraph
{Similar families exist for cubical meshes rather than simplicial triangulations.}
\end{tcolorbox}
\par\endgroup
% Original PDF 106 is a blank production leaf with printer marks, retained in the source.
